\section{Cold Hardiness Data}
\label{sec:data}

This study used the cold hardiness of endo\textendash and ecodormant primary buds from about 30 diverse field-grown grapevine cultivars measured since 1988 at the WSU Irrigated Agriculture Research and Extension Center (IAREC). Locations include the vineyards at IAREC, Prosser, WA (46.29°N lat., -119.74°W long.), the WSU-Roza Research Farm, Prosser, WA (46.25°N lat., -119.73°W long.), and the Ste. Michelle Wine Estates, Paterson, WA (45.96°N lat.; -119.61°W long.). Cane samples containing dormant buds were collected daily, weekly, or at 2-week intervals from leaf fall in October to bud swell in April---i.e., dormant season. The collected samples were analyzed using differential thermal analysis (DTA) \cite{mills_cold-hardiness_2006}. DTA requires putting the samples in a thermoelectric module that senses low-temperature exotherms (LTEs) resulting from the freezing events of individual buds. The module is placed in a controlled chamber, where LTEs  are monitored and registered as the temperature decreases. The result is a measurement of the lethal temperatures at which 10\%, 50\%, and 90\% of the bud population die, denoted by $LTE_{10}$, $LTE_{50}$, and $LTE_{90}$, respectively. Additionally, daily environmental data (e.g., max and min air temperature) from the closest on-site weather station to each vineyard was obtained using the API provided by AgWeatherNet \cite{AgWeatherNet}. %The three stations used are Prosser.NE (46.25°N latitude, -119.74°W longitude), Roza.2 (46.25°N latitude, -119.73°W longitude), and Paterson.E (45.94°N latitude, -119.49°W longitude).
%In addition to air temperature, the AgWeatherNet stations record a host of numerous other weather variables including but not limited to: humidity, solar radiation, and wind speed.

Thus, for each cultivar, there is a temporal dataset with a varying number of seasons containing daily weather data along with cold hardiness LTE values for the days that samples were collected. For the purpose of this study, a \emph{season} is said to span from September 7th to May 15th---a conservative interval containing the full dormancy period \cite{ferguson_dynamic_2011, ferguson_modeling_2014}. 

\paragraph{Data Summary.} Table \ref{tab:data-description} shows a summary of the number of years of data collected for the five different grape cultivars that have been selected for this study based on their market importance ($n=3,629$ samples). 
%The interval of years in parenthesis represents the in-between seasons for which there was no LTE data collected---e.g., (1990-1992) means seasons 1990-1991 and 1991-1992 have no LTE data.
Each cultivar dataset contains a row for each day of data collection. %Note, since cold hardiness was not measured on each day of a season, some rows  do not contain LTE data. 
%Below we list selected fields of the data contained in each row.
The key (selected) temporal fields include:
\begin{itemize} 
\item DATE: The date of  observation. 
%\item SEASON: Years comprising the dormant season from September 7th (JDAY 250) to May 15th (JDAY 500).
\item SEASON\_JDAY: Julian day is an integer that represents the count of the number of days since the beginning of the year. For the dormant season, the JDAY continues at the new year---i.e., starts on JDAY 250 (Sep 7th) and ends with JDAY 500 (May 15th). 
%in a typical year, Dec 31 is JDAY 365, then Jan 1 is JDAY 366, and so on. The dormant season typically 
%\item AWN\_STATION: The closest AgWeatherNet station from where the environmental readings are taken.
%\item DORMANT\_SEASON: It will have a value of 1 if the date belongs to the dormant season or 0 if it does not.
% \item LTE10 (when available): Lethal temperature observed from vineyard samples when 10\% of the buds die. In degrees Celsius. 
% \item LTE50 (when available): Lethal temperature observed from vineyard samples when 50\% of the buds die. In degrees Celsius.
% \item LTE90 (when available): Lethal temperature observed from vineyard samples when 90\% of the buds die. In degrees Celsius.
\begin{table}[th!]
\centering
\resizebox{1\columnwidth}{!}{
\begin{tabular}{ |l|l|c|c| }
\hline
\multicolumn{1}{|c|}{\textbf{Cultivar}}             &\multicolumn{1}{|c|}{\textbf{LTE Data Seasons}}                      &\makecell{\textbf{Years of} \\ \textbf{LTE Data}}    &\makecell{\textbf{LTE Total} \\ \textbf{Samples}} \\\hline
Cabernet Sauvignon (CS)	 &1988-2022	                                    &34                         &829 \\\hline
Chardonnay (CH)	         &1996-2022	                                    &26                         &783 \\\hline
%Concord (CD)	             &1988-2022 (1990-1992, 1993-1997, 1998-1999)	&27                         &484 &Juice\\\hline
Concord (CD)	             &1988-'90,'92-'93,'97-'98,'99-2022	&27                         &484 \\\hline
Merlot (MR)	             &1996-2022	                                    &26                         &897 \\\hline
Riesling (WR)             &1988-2022	                                    &34                         &636 \\\hline
\end{tabular}
}
\caption{Summary of LTE data for selected grape cultivars.}
\label{tab:data-description}
\end{table}
\item LTE values at $LTE_{10}$, $LTE_{50}$, and $LTE_{90}$: in $^{\circ}$C. 
%\item $LTE_{50}$ (when available): in degrees Celsius.
%\item $LTE_{90}$ (when available): in degrees Celsius.
\item MIN\_AT, AVG\_AT, MAX\_AT: Minimum, average, and maximum air temperatures respectively, as observed at 1.5 meters above the ground (in degrees Celsius). 
\end{itemize}
